\documentclass[main]{subfiles}


\title[AutoML: Introduction]{AutoML Motivation}
\subtitle{How to Define AutoML Formally?}
\author[Marius Lindauer]{Marius Lindauer}
\date{}

\titlebackground*{images/AutoML_AutoDL}

%\institute{}
%\date{}

\begin{document}
	
	\maketitle
	

%----------------------------------------------------------------------
%----------------------------------------------------------------------
\begin{frame}[c]{Machine Learning Task}

Let  

\begin{description}
    \item[$\dataset$] be a dataset, typically in the form of $\dataset = (\xI{i},\yI{i})_{i \in 1 \ldots N}$ with $N$ data points,
    \item[$\Hspace$] be a hypothesis space with all possible predictive models $f$ under consideration,
\end{description}

the machine learning task is to learn a predictive model $\fh \in \Hspace$\\ that minimizes the empirical risk of errors:

$$ \fh \in \argmin_{f \in \Hspace} \riske(f , \dataset) $$

\end{frame}
%-----------------------------------------------------------------------
%----------------------------------------------------------------------
\begin{frame}[c]{Empirical Risk: Assessing How Well a Model Performs?}

Let  

\begin{description}
    \item[$\dataset$] be a dataset, typically in the form of $\dataset = (\xI{i},\yI{i})_{i \in 1 \ldots N}$ with $N$ data points,
    \item[$\loss$] be a loss function on two labels.
\end{description}

The empirical risk of a model $f$ is defined by:

$$\riske(f) := \sum_{i=1}^{N} \loss\left(\yI{i}, f(\xI{i})\right)$$

\bigskip 

Note: Some papers also use $\mathcal{L}$ to denote the (empirical) loss.

\end{frame}
%-----------------------------------------------------------------------
%----------------------------------------------------------------------
\begin{frame}[c]{How to Get the Model?}

Let  

\begin{description}
    \item[$\dataset$] be a dataset, typically in the form of $\dataset = (\xI{i},\yI{i})_{i \in 1 \ldots N}$ with $N$ data points,
    \item[$\pcs$] be a configuration space of all considered design decisions\footnote{On the previous slides, we implicitly assumed a given $\conf$. By considering $\conf$ explicitly, we open up the hypothesis space.}, 
    \item[$\Hspace$] be a hypothesis space with all possible predictive models $f$ under consideration, 
\end{description}

such that $\mathcal{I}$, being a learning algorithm (a.k.a. inducer), infers the model from the data and a configuration $\conf \in \pcs$:

\vspace*{-1em}
\begin{eqnarray}
\mathcal{I}: \datasets \times \pcs \to \Hspace \qquad (\dataset, \conf) \mapsto f\\
\text{s.t. } \fh \in \argmin_{f \in \Hspace} \riske(f,  \dataset)
\end{eqnarray}

\end{frame}
%-----------------------------------------------------------------------
%----------------------------------------------------------------------
\begin{frame}[c]{What is $\conf \in \pcs$?}

Possible elements described via $\pcs$ could include:
 
\begin{description}
    \item[Hyperparameter Configuration] For example, $\conf$ could specify the setting of a learning rate and data augmentation strength.
    \item[Neural Architecture] For example, $\conf$ could specify the operators and connections of a cell in a neural network.
    \item[Pipeline] For example, $\conf$ could specify that we first run standardization and then a PCA before feeding it into a linear model.
    \item[Ensemble] For example, $\conf$ could specify which models to combine in an ensemble.
\end{description}

\end{frame}
%-----------------------------------------------------------------------
%----------------------------------------------------------------------
\begin{frame}[c]{\alert{Automated} Machine Learning Problem}

Let  

\begin{description}
    \item[$\datasettrain$] be a training dataset,
    \item[$\datasetval$] be a validation dataset\footnote{typically disjoint from $\datasettrain$. In principle any resampling strategy could be used.}, and
    \item[$\pcs$] be a configuration space of all considered design decisions (e.g., hyperparameters or architectural decisions),
    \item[$\riskg$] be a generalization risk\footnote{The risk or loss metrics on the training and validation dataset do not have to be the same}.
\end{description}

The automated machine learning task is to find a configuration $\optconf$ of the inducer $\mathcal{I}$ such that the risk on the validation data is minimized

$$\optconf \in \argmin_{\conf \in \pcs} \riskg(\mathcal{I}(\datasettrain, \conf), \datasetval)$$

\end{frame}
%-----------------------------------------------------------------------
%----------------------------------------------------------------------
\begin{frame}{Summary} 

Most important insights from this video:

\begin{enumerate}
    \item The learning of the machine learning model focuses on minimizing the training error, which should eventually lead to better generalization error overall
    \item An inducer (a.k.a. learning algorithm) and the configuration induce a model from a given dataset
    \item Automated machine learning finds the optimal configuration of the inducer to minimize the generalization error
    \begin{itemize}
        \item \alert{Attention}: This requires special validation and evaluation schemes (discussed in other sessions)
        \item \alert{Note}: The configuration space is w.r.t. the inducer in our definition; it can cover the entire data science pipeline or even code generation
    \end{itemize}
\end{enumerate}

\end{frame}

\end{document}
